% ----------------------------------------------------------
% Capítulo 1 — Introdução
% ----------------------------------------------------------
\chapter{Introdução}
\label{cap:introducao}

A modelagem da volatilidade de ativos financeiros constitui um dos pilares
da econometria financeira moderna, com aplicações diretas na gestão de
risco, na precificação de derivativos e na alocação de
portfólios \cite{tsay2010}. Desde a introdução dos modelos ARCH por
\citeonline{engle1982} e sua generalização GARCH por
\citeonline{bollerslev1986}, a literatura tem avançado significativamente
na compreensão da dinâmica temporal da variância condicional de retornos
financeiros. Entretanto, a qualidade dos modelos estimados depende
fundamentalmente da forma com que os dados brutos de mercado são
pré-processados e agregados --- uma etapa frequentemente negligenciada na
prática.

A forma mais comum de agregação de dados financeiros de alta frequência é
a construção de barras temporais (\emph{time bars}), nas quais as
observações são agrupadas em intervalos fixos de tempo. Contudo,
\citeonline{lopez2018advances} demonstrou que essa abordagem introduz
distorções estatísticas relevantes: sobre-amostragem em períodos de baixa
atividade, sub-amostragem em períodos de alta atividade, e retornos com
propriedades que se afastam das premissas de modelos econométricos
tradicionais. Como alternativa, o autor propôs as \emph{dollar bars},
barras baseadas no volume financeiro transacionado, que produzem séries com
propriedades estatísticas mais favoráveis à modelagem.

A questão que surge naturalmente é: \emph{em que medida a escolha do método
de amostragem --- \emph{time bars} versus \emph{dollar bars} --- afeta a
qualidade das estimativas de volatilidade obtidas por modelos
ARMA-GARCH?} Esta pergunta tem relevância prática direta, pois a
estimativa da variância condicional é insumo para decisões de investimento,
\emph{hedging} e controle de risco.

Adicionalmente, a avaliação de modelos em séries temporais financeiras
requer cuidados metodológicos específicos. A dependência temporal entre
observações invalida métodos tradicionais de validação cruzada, e a prática
de selecionar o ``melhor modelo'' com base em métricas \emph{in-sample}
pode conduzir a sobreajuste (\emph{overfitting}). O arcabouço de Validação
Cruzada Combinatória Purgada (CPCV), proposto por
\citeonline{lopez2018advances}, endereça essas limitações ao gerar
múltiplos caminhos de \emph{backtest} e permitir a estimação das estatísticas em séries
temporais com maior rigor.

% ----------------------------------------------------------
\section{Objetivos}
\label{sec:objetivos}
% ----------------------------------------------------------

\subsection{Objetivo Geral}
\label{subsec:objetivo_geral}

Analisar o impacto da escolha do método de amostragem de dados financeiros
--- \emph{time bars} versus \emph{dollar bars} --- sobre as propriedades
estatísticas das séries resultantes e sobre a qualidade preditiva de modelos
ARMA-GARCH para a variância condicional de retornos, utilizando o
estimador Two-Scale Realized Volatility (TSRV) para a variância realizada,
como referência e o arcabouço CPCV para avaliação comparativa robusta.

\subsection{Objetivos Específicos}
\label{subsec:objetivos_especificos}

\begin{enumerate}
    \item Construir séries de retornos a partir de \emph{time bars} e
          \emph{dollar bars} para um ativo financeiro selecionado, comparando
          suas propriedades estatísticas descritivas;
    \item Estimar a variância realizada por meio do estimador TSRV;
    \item Ajustar modelos ARMA-GARCH conjuntos sobre ambos os tipos de
          barra, selecionando as ordens ótimas por critérios de informação;
    \item Avaliar e comparar o desempenho preditivo dos modelos utilizando métricas de erro de previsão de
          volatilidade (MSE, MAE, QLIKE, MZ).
    \item Comparar o desempenho preditivo dos modelos com o estimador TSRV.
\end{enumerate}

% ----------------------------------------------------------
\section{Justificativa}
\label{sec:justificativa}
% ----------------------------------------------------------

A escolha do método de amostragem de dados financeiros é uma decisão
metodológica de primeira ordem que afeta toda a cadeia subsequente de
modelagem e inferência. Apesar de amplamente discutida na literatura de
aprendizado de máquina financeiro \cite{lopez2018advances}, a interação
entre métodos de amostragem alternativos e modelos econométricos clássicos
de volatilidade --- particularmente a família GARCH --- permanece pouco
explorada na literatura acadêmica brasileira.

Este trabalho justifica-se pela necessidade de:

\begin{enumerate}
    \item \textbf{Preencher uma lacuna metodológica}: a maioria dos estudos
          que empregam modelos GARCH utiliza \emph{time bars} sem questionar
          as implicações dessa escolha. Uma análise comparativa rigorosa
          contribui para a compreensão de como o método de amostragem afeta
          as estimativas de volatilidade;
    \item \textbf{Adotar uma avaliação robusta}: o uso do CPCV representa uma abordagem metodologicamente
          superior ao \emph{walk-forward} simples, comumente empregado em
          trabalhos aplicados;
    \item \textbf{Integrar estimadores eficientes}: o estimador TSRV é uma alternativa de estimação em cima de tick data para a estimação de variância realizada mais precisa do que o retorno quadrático simples, permitindo uma avaliação mais fidedigna dos modelos.
\end{enumerate}
